\section{Aims of dissertation thesis}
\label{sec:aims}

The previous chapters established the technological foundations of blockchain and identified the Internet of Value as an emerging paradigm that enables direct, peer-to-peer exchange of assets through decentralized infrastructure. Our analysis of the state of the art revealed three critical gaps that prevent the full realization of this vision.

First, despite the growing diversity of digital assets, there is no standardized framework for multi-dimensional, goal-oriented asset valuation. Current approaches rely on market price alone and fail to capture the broader value dimensions, such as utility, durability, and scarcity, that are essential in the Internet of Value \cite{Lohmer2024}.

Second, existing blockchain account models, whether externally owned accounts or smart contract wallets, lack the abstractions required for autonomous agent operation. There are no specialized primitives for policy enforcement, risk isolation, and auditable decision-making that remain compatible with existing DeFi protocols.

Third, governance processes in decentralized ecosystems rely heavily on manual evaluation by community members. As the complexity and volume of treasury proposals grow, this approach becomes infeasible and prone to bias \cite{Valastin2024proto}.

With the maturity of large language models and emerging technologies such as the Model Context Protocol (MCP), which ensures reproducible, deterministic tool invocation, we can now delegate complex tasks to AI models capable of autonomous execution \cite{Sumers2024, Wu2023}. The intersection of these agentic systems with blockchain creates an opportunity: while blockchain serves as the immutable settlement layer, intelligent agents provide the reasoning capability to support better decision-making. Instead of requiring users to navigate complex interfaces and protocol-specific constraints, we can leverage AI's flexibility to interpret human goals into executable blockchain transactions.

The research we are proposing represents a new approach to the Internet of Value, a paradigm shift in the way we use digital assets over blockchain technology using agentic systems. The main focus areas are captured in three research questions:

\begin{enumerate}
   \item How can multidimensional asset value be formally defined and evaluated by autonomous agents within the Internet of Value?
   \item Propose the new architecture for an agetic system in the context of blockchain and the Internet of Value.
   \item How can autonomous agents leverage contextual value assessment to support objective decision-making in on-chain treasury governance?
\end{enumerate}

The first research question addresses the challenge of effective asset operation within the Internet of Value. Value is fragmented across multiple ecosystems, each with distinct token standards, governance mechanisms, and consensus rules. Evaluating assets requires navigating various block explorer interfaces, understanding protocol-specific caveats, and aggregating data from disparate marketplace APIs. We aim to propose a formal value framework for measuring multidimensional asset value that provides objective criteria for evaluation while remaining flexible enough to adjust weights based on operational goals. Additionally, we aim to develop on-chain execution primitives that enable autonomous agents to own, operate, and manage assets with appropriate policy constraints and user sovereignty.

The second research question focuses on developing a novel architecture that bridges human intent with on-chain execution.
The standard approach to building an agentic solution is tight coupling of components, which ensures a somewhat straightforward implementation. Therefore, agents are usually built with a specific purpose and can process only one task\footnote{Just sending tokens on a particular chain}.
We explored an operational framework in which agents are divided into multiple independent layers, acting as a flexible pipeline rather than a transactional facilitator. One of the critical layers in this framework is an agent with contextual asset intelligence. AI agents can understand not only the technical properties of blockchain assets but also their contextual significance and relationship networks within the multi-dimensional value space. This fundamentally shifts the view of asset valuation from metric evaluation towards a deeper understanding of value representation across multiple dimensions.

Using a multi-layered approach, we can address significant asset challenges in decentralized environments while maintaining flexibility and enabling composability. The main benefit is that we can stay in the loop throughout the process and are not left out of the decision-making. Our architecture consists of specialized layers. First is valuation agents for data acquisition and metric computation across the three value dimensions. The second one is decision agents for intent-based contextualization via weighted evaluation. Most importantly, the execution agents for on-chain operationalization via NFTAA primitives.

The third research question focuses on the application of agentic systems in the domain of on-chain governance. As the Internet of Value expands across multiple blockchains and incorporates an ever-growing diversity of asset types, the challenge of funding allocation becomes increasingly complex. Projects seeking funding through quadratic funding models or on-chain treasury proposals require evaluation of extensive on-chain and off-chain data, a task that exceeds the capacity of human evaluators. We aim to demonstrate that our value framework can be adapted to the governance domain, enabling agents to produce objectively generated, unbiased assessments based on resource scarcity, long-term sustainability, and community benefit. This is particularly relevant for systems like Polkadot, where the treasury holds substantial resources that must be allocated efficiently and transparently.

Along with the main research questions, we have identified specific challenges that correspond to each research area:

\begin{enumerate}
  \item There is no framework that can objectively evaluate heterogeneous asset types while adapting to diverse user intents.
  \item Current blockchain systems lack a standard architecture for integrating autonomous agents with on-chain execution in a modular, extensible manner.
  \item Manual evaluation of on-chain treasury proposals is biased and does not scale, resulting in inconsistent decision-making.
\end{enumerate}

The following chapter presents our proposed design that addresses these research questions and challenges through a layered architecture combining formal value assessment, specialized agentic components, and novel on-chain execution primitives.
